\chapter{The Governance Contest: Rules, Standards, and Who Writes Them}
\chaptermark{The Governance Contest}
\label{ch:governance}

Technology wars are settled twice: once in the factories, and once in the rulebooks. Part~I's precedent industries make the point. Solar was decided by cost curves, but the residue of the war is a standards regime---efficiency floors, grid codes, module certifications---written substantially in China. Batteries were decided by chemistry and scale, and the aftermath is a Chinese export-control regime on cathode process technology. Automotive safety and charging standards followed the same path. In each case the winner of the industrial contest inherited the pen.

AI's rulebook is being written now, and this chapter argues that it is being written in four separate arenas by three actors who are not, for the most part, competing with one another---each has occupied the arena its own institutional strengths suit, and the arena that matters most is occupied by nobody (Figure~\ref{fig:gov}). Understanding that asymmetry is necessary for the strategy chapters that follow, because the standard Western framing---``the US regulates lightly, the EU regulates heavily, China regulates authoritarianly''---misdescribes the contest badly enough to lose it.

\begin{figure}[htbp]
\centering
\begin{tikzpicture}[
  every node/.style={font=\scriptsize},
  arena/.style={draw=inkMut, line width=0.6pt, rounded corners=2pt, fill=surf,
                align=center, inner sep=4pt, text width=26mm, minimum height=9mm},
  actor/.style={align=center, font=\scriptsize\color{inkSec}, text width=25mm},
  lead/.style={draw=none, fill=sOne!18, rounded corners=2pt, align=center,
               inner sep=3pt, text width=25mm, minimum height=8mm},
  weak/.style={draw=none, fill=surf, rounded corners=2pt, align=center,
               inner sep=3pt, text width=25mm, minimum height=8mm},
  none/.style={draw=gridLn, line width=0.5pt, densely dashed, fill=white,
               rounded corners=2pt, align=center, inner sep=3pt, text width=25mm, minimum height=8mm},
]
\node[actor] at (0,1.1) {};
\node[actor] at (3.0,1.1) {\textbf{United States}};
\node[actor] at (6.0,1.1) {\textbf{European Union}};
\node[actor] at (9.0,1.1) {\textbf{China}};
\node[actor] at (12.0,1.1) {\textbf{Neutral / unowned}};

\node[arena] at (0,0.3) {\textbf{Model \& content}};
\node[weak] at (3.0,0.3) {sectoral, EO churn};
\node[lead] at (6.0,0.3) {AI Act: GPAI duties};
\node[weak] at (9.0,0.3) {algorithm \& genAI filings};
\node[none] at (12.0,0.3) {---};

\node[arena] at (0,-1.0) {\textbf{Compute \& export}};
\node[lead] at (3.0,-1.0) {BIS controls, entity lists};
\node[weak] at (6.0,-1.0) {follows US, partially};
\node[weak] at (9.0,-1.0) {counter-controls: minerals, battery tech};
\node[none] at (12.0,-1.0) {---};

\node[arena] at (0,-2.3) {\textbf{Standards}};
\node[weak] at (3.0,-2.3) {firm-led, consortium};
\node[weak] at (6.0,-2.3) {process-led};
\node[lead] at (9.0,-2.3) {ISA, interconnect, GB standards, WAICO};
\node[none] at (12.0,-2.3) {---};

\node[arena] at (0,-3.6) {\textbf{Trust \& certification}};
\node[weak] at (3.0,-3.6) {voluntary, lab-led};
\node[weak] at (6.0,-3.6) {conformity assessment (paper)};
\node[weak] at (9.0,-3.6) {domestic filing regime};
\node[lead, fill=sTwo!22] at (12.0,-3.6) {\textbf{vacant}};

\node[font=\scriptsize\color{inkSec}, align=left, anchor=north west, text width=132mm]
  at (-1.35,-4.4) {Shaded cells mark where an actor currently sets the operative rules. Each governs the arena it is
  strongest in---the US by denial, the EU by regulation, China by standards and institution-building---and no
  actor governs the arena that the argument of Chapters~\ref{ch:trust} and~\ref{ch:order66} shows matters most.};
\end{tikzpicture}
\caption[The governance arena map]{Four arenas, three philosophies, one vacancy. Governance of AI is not a single
contest but four, and the actors are not competing in the same one: each has occupied the arena its institutional
strengths suit. The trust-and-certification arena---provenance, artifact lineage, integration-level evaluation---is
the one whose absence Chapter~\ref{ch:order66} shows to be systemically consequential, and it is unowned.}
\label{fig:gov}
\end{figure}


\section{Three governing philosophies}

\textbf{The United States governs by denial.} Its principal instrument of global AI governance is not a statute but an export-control regime: the October 2022 and December 2024 rules, the entity lists, the HBM and equipment controls, the short-lived AI Diffusion Rule, and the case-by-case licensing of accelerators to named countries. This is genuinely extraterritorial regulation---the foreign direct product rule reaches any fab using US tooling---and for a period it functioned as the world's operative AI governance, because access to compute gated everything else. Domestic model governance, by contrast, has been sectoral, voluntary, and subject to executive-order churn; the US has no general AI statute. The strategic weakness of governance-by-denial is the one this book has documented at length: it produces substitution rather than compliance, and each denied layer becomes a subsidized domestic industry (Chapter~\ref{ch:playbook}). Its second weakness is that denial confers no positive standard---no definitions, no conformity procedure, no registry---so when the denial fails there is nothing left behind. The pattern held to form in August 2026: the White House finalized a voluntary frontier-model evaluation framework in a closed-door session with the major US laboratories---and declined to publish the document itself~\cite{wh-framework}. Voluntary and opaque where the Chinese regime is mandatory and registered: whatever one thinks of either pole, only one of them is accumulating the definitions, filings, and institutional muscle that outlast a political cycle.

\textbf{The European Union governs by regulation.} The AI Act is the world's most complete general AI statute: risk tiers, obligations for general-purpose model providers, transparency and documentation duties, systemic-risk designation above a compute threshold, and partial exemptions for open-source releases. It is real, it is extraterritorial in effect through the market it controls, and it has produced the fullest public vocabulary yet for talking about model obligations. Its weakness is the one Chapter~\ref{ch:strategies} named: regulation is not production. A jurisdiction that neither trains the frontier models nor manufactures the accelerators is writing rules whose subjects are elsewhere, and its leverage decays with its share of the deployment market. The open-source exemptions are also where the Act is least settled, because ``open'' in the Act's sense and ``open weights'' in the industry's sense are not the same object---a gap the next section takes up.

\textbf{China governs by standards and institution-building.} This is the arena Western commentary consistently under-weights, and it is where the pattern of Part~I repeats most exactly. Domestically, China has a functioning filing-and-review regime for algorithms and generative services---the earliest such regime anywhere---but the strategically consequential activity is external and infrastructural: the mandated unified AI-chip interconnection ecosystem; national standards (GB-series) that become procurement conditions; the open-sourcing of interconnect specifications and software stacks so that the de facto standard is Chinese by adoption rather than by decree (Chapter~\ref{ch:compute}); the largest single bloc in the open-ISA consortium; and, at the international level, the Global AI Governance Action Plan, the UN capacity-building group co-chaired with a developing-country partner across some eighty states, BRICS AI centers, and the World AI Cooperation Organization headquartered in Shanghai with twenty-nine founding members, none of them a major Western democracy~\cite{action-plan,carnegie-pivot,waico}. \emph{This is what winning a standards war looks like from the inside}: not a treaty imposed, but a stack given away until it is the default, and then a forum built around the default.

\section{Standards are the load-bearing arena}

The technical standards layer is where governance becomes irreversible, because a ratified standard survives the political cycle that produced it and because compliance is voluntary in name and compulsory in practice. Four bodies matter more than any ministry.

\emph{Instruction sets.} RISC-V International is Swiss-incorporated specifically to be sanction-proof, and twelve of its twenty-four premier members are Chinese~\cite{csis-riscv}. The matrix-extension task groups now standardizing the operations that AI silicon executes are the single most consequential unglamorous venue in this book, and Congressional letters urging restrictions on US participation have produced no action for the reason Commerce itself conceded: restricting participation in an open standard harms the restricting party (Chapter~\ref{ch:compute}). \emph{Memory.} JEDEC sets LPDDR6 and SOCAMM2---the standards on which the capacity-first architecture of Section~\ref{sec:lpddr6} depends, and on which a Chinese producer is now at the leading edge of implementation. \emph{Interconnect.} The contest between an open, state-adopted specification and a Western consortium that excludes Chinese membership is a live test of whether openness or exclusivity wins standards adoption; the historical record of this book suggests one answer. \emph{Model and system standards.} ISO/IEC JTC1 SC42, IEEE working groups, and ITU study groups are where terminology, evaluation methods, and management-system requirements are fixed, and where participation intensity translates directly into definitional power years later.

``Standard,'' moreover, names four different degrees of reality, and the maturity map used here earns its place because it converts standards talk into standards accounting. A specification climbs: (1) \emph{de jure specification}---ratified text; (2) \emph{open reference implementation}---code or silicon anyone can run; (3) \emph{conformance suite}---a test that adjudicates claims of compliance; (4) \emph{installed-base default}---what new deployments assume without deciding. Rent lives at stage 4; committees fight over stage 1; and the distance between them is where standards strategies succeed or die. Applying the map to this book's contested standards: RISC-V RVV~1.0 sits at stage 3--4 (ratified, multiple implementations, conformance in place, default in Chinese silicon policy); the IME/AME matrix extensions at stage 1-in-progress with stage 2 explicitly scheduled (Section~\ref{sec:riscv}); UnifiedBus at stage 2--3 inside China (published spec, deployed SuperPoDs, mandated adoption) and stage 0 outside it; UALink and Ultra Ethernet at stage 1--2 with the installed base still NVLink's; CXL ratified for years yet still climbing toward stage 4 against socket economics; LPDDR6 and SOCAMM2 at stage 4 (JEDEC-ratified, mass-produced, default in new server designs); ONNX at a cautionary stage 3-without-4 (conformant, implemented, and routinely bypassed by direct framework-to-kernel paths); OpenXLA at stage 2 betting on 4; model-license terms at stage 4 by sheer adoption with no conformance machinery at all; artifact provenance (signed manifests, model BoMs) at stage 1 at best---the certification vacancy of Chapter~\ref{ch:trust} restated as a maturity score; and MLPerf, not formally a standard, operating the field's only functioning stage-3 conformance culture. The strategic reading of the map: China's distinctive play is to \emph{skip stage politics}---ship stages 2 and 4 (reference implementation, mandated installed base) and let stage 1 ratify what already exists; the West's distinctive failure is stranding specifications at stage 1--2 while defending installed bases it assumes are permanent.

The lesson from Part~I is that standards leadership \emph{tends} to follow manufacturing leadership with a lag of five to ten years, and then locks it in for twenty. The qualifier matters, and this book should not overclaim it: industrial scale confers influence, not automatic control. International standards bodies have voting rules, patent policies, national delegations, and safety and certification constituencies that do not simply ratify whoever ships most units. The counter-cases are instructive. China dominates solar manufacturing but the IEC's PV standards machinery remains substantially European- and Japanese-shaped. It leads battery production yet automotive functional-safety and charging-interface standards remain contested, with CCS, CHAdeMO, NACS, and GB/T coexisting rather than converging on the largest producer's specification. Its 5G patent position translated into real standards-essential influence---but only after a decade of sustained participation, and it did not carry the adjacent semiconductor standards with it. Conversely, China has converted scale into de facto standardization fastest where it controlled \emph{both} the specification and the reference implementation and gave them away---which is precisely the RISC-V and open-interconnect playbook, and precisely why that arena, rather than the formal committees, is the one this chapter treats as load-bearing. The West is currently spending its standards influence defending chokepoints in arenas where it is strong and neglecting arenas where the terms of the next decade are being written.

\section{The open-weight governance problem}

Open weights have broken the categories every existing regime relies on, and no jurisdiction has repaired them.

\emph{What is ``open''?} Most Chinese frontier releases ship under MIT or Apache~2.0---licenses that are unambiguously open for the \emph{artifact} while disclosing little about training data, recipes, or evaluation. The industry's convention (``open source model'') overstates what is released; this book uses \emph{open-weight} throughout for that reason. The definitional fight is not pedantic: the EU AI Act's exemptions, the OSI's attempt at an open-source-AI definition, and emerging model-distribution licenses each draw the boundary differently, and where the boundary falls determines which obligations attach to the models that a majority of the world's developers now use.

\emph{Who is the provider?} Regulatory regimes assume a provider who places a system on the market and remains accountable for it. A quantized derivative of a merged fine-tune of an open checkpoint, redistributed by an individual and deployed on-premises by a third party, has no such provider. Chapter~\ref{ch:order66} showed that this chain is also the security-critical one; here the point is that it is a governance vacuum, not merely an engineering one.

\emph{Do compute thresholds still work?} The dominant technical instrument in both the EU's systemic-risk designation and US executive practice has been a training-compute threshold---a sensible proxy while frontier training required a hyperscale campus. The architecture argument of Chapter~\ref{ch:compute} attacks it directly: if capacity-first memory and sparse recipes bring frontier-class pretraining within reach of 128--256-socket sovereign clusters drawing 0.5--1\,MW (Section~\ref{sec:lpddr6}), then threshold-and-registry governance faces many more, much smaller, much less visible training runs, most of them outside any single jurisdiction. Algorithmic efficiency compounds the problem from the other side: a fixed FLOP threshold denotes a falling capability level every year. \emph{Decentralization is not only an industrial forecast; it is a governance forecast, and the instruments in force today do not survive it.}

\section{Dual-use science: the governance failure mode that would hurt most}

There is a specific way governance could go badly wrong for the readership of this book, and it deserves naming. The scientific domains where AI is most valuable---biology, chemistry, materials, robotics---are inherently dual-use. A broad restriction written to prevent misuse of capable models can, without any malice, foreclose legitimate research: capability evaluations that gate model access, data-sharing restrictions that break international collaboration, or liability regimes that make hosting a scientific model at a university uninsurable~\cite{consortium-record}.

Two structural observations follow. First, institutions that hold their own models, their own evaluation evidence, and their own provenance records are in a far stronger position to argue for proportionate treatment than institutions that can only relay a vendor's assurances---which is one of the three risks that makes open pretraining an insurance policy rather than a vanity project (Chapter~\ref{ch:consortium}). Second, the same evidence that would support proportionate governance is the evidence the certification arena would produce if anyone were producing it. Scientific self-governance in this domain is not a way of avoiding rules; it is the only realistic way of getting rules that a working scientist can live with.

\section{The vacant arena, and why it will be filled}

Figure~\ref{fig:gov}'s fourth row is the argument of this chapter. Nobody governs trust and certification for the open commons: signed artifact manifests, reproducible lineage across derivative chains, integration-level injection and backdoor evaluation, published behavioral and censorship audits, national or international mirrors. The US regime cannot fill it, because governance-by-denial produces no positive procedure and because it will not certify artifacts it wishes did not exist. The EU regime has the paper apparatus---conformity assessment, technical documentation, notified bodies---but not the technical operations, and its subjects are largely outside its market. China's forum-building is an explicit bid for the role, and it is the only bid currently on the table; a certification authority operated by the commons' principal supplier is not obviously worse than none, but it is not what a buyer would design.

The arena will be filled, because the deployment layer cannot scale without it. Regulated sectors, safety-critical industry, defense-adjacent research, and any government procuring at scale will all require what certification provides, and they will accept it from whoever offers it credibly first. That is the same dynamic by which certification stalled the C919 and the sintilimab application (Chapter~\ref{ch:analogy})---and, read in the other direction, the same dynamic that would let a credible neutral institution acquire durable influence over a commons it did not build. Chapter~\ref{ch:consortium} argues that this is the highest-leverage institutional position available to any country that is not a superpower, and Chapter~\ref{ch:sizing} supplies the demand-side numbers that make the associated infrastructure affordable if pooled.

\section{Three governance futures}

\emph{Convergence} (least likely, most desirable): shared definitions, mutual recognition of conformity assessments, an internationally operated evaluation and provenance layer, and threshold instruments replaced by capability-and-deployment-based ones. This requires the two principals to accept a referee neither controls---which has happened before in aviation and pharmaceuticals, and only after accidents.

\emph{Bifurcation} (most likely): two standard stacks with partial technical interoperability and no mutual recognition---a Western regulatory bloc governing its own market by statute, a Chinese-led standards-and-forum bloc governing the majority of the world's deployments by adoption, and third countries certifying twice or choosing. This is Chapter~\ref{ch:trade}'s Scenario~III expressed in institutions rather than in market share, and it is where present trends point.

\emph{Fragmentation} (the failure case): compute thresholds obsoleted by decentralized training, open-weight provider obligations unenforceable against derivative chains, no certification layer anywhere, and governance retreating to national procurement rules and incident response after the fact. This is not a stable equilibrium; it is the state the system occupies until a sufficiently public failure---of the compound-agentic kind Chapter~\ref{ch:order66} documents---forces one of the other two.

The strategic conclusion for the chapters that follow is narrow and firm. The rules of the AI era are not being written where most Western attention is directed. They are being written in instruction-set task groups, memory standards committees, interconnect specifications, license texts, and a new international organization in Shanghai---and in one arena that is still empty, whose occupant will hold, over the deployment layer, the leverage that export controls were supposed to provide over the supply layer.
